\pagebreak
\section{Zadaci}
%TODO: DODATI ZADATAK ZA RAZLICITE STRUJE POLARIZACIJE

\newtheorem{zadatak}{Zadatak}

\begin{zadatak}
    Sastaviti kolo neinvertujućeg pojačavača, koristiti operacioni pojačavač AD8040, gde su $\pm V_{cc}=\pm10\,V$, a $R_1 = 100 \,k\Omega$ i $R_2 = 1 \,M\Omega$, i odrediti:

    \begin{enumerate}
        \item [a)] Očekivani izlaz pojačavača, ako se na ulaz dovodi signal $V_{in} = 0.9\,V sin(20\,\pi t)$,
        \item [b)] Vrednost struja polarizacije za kolo,
        \item [c)] Prenosnu karakteristiku kola,
        \item [d)] Kompenzaciju uticaja struja polarizacije na izlaz kola,
        \item [e)] Izlaz kola nakon kompenzacije za ulaz iz a).
    \end{enumerate}
\end{zadatak}

\begin{zadatak}
    Sastaviti kolo invertujućeg pojačavača, koristiti operacioni pojačavač ADA4891, gde su $\pm V_{cc}=\pm10\,V$, a $R_1 = 100 \,k\Omega$ i $R_2 = 1 \,M\Omega$, i odrediti:
    
    \begin{enumerate}
        \item [a)] Očekivani izlaz kola ako se na ulaz dovede signal $V_{in} = 0.1\,V\,sin(20\,\pi t)$,
        \item [b)] Vrednost naponskog ofseta u kolu,
        \item [c)] Prenosnu karakteristiku pojačavača,
        \item [d)] Vrednost otpornika za kompenzaciju naponskog ofseta i
        \item [e)] Vrednost izlaznog signala za ulazni signal iz a).
    \end{enumerate}
\end{zadatak}

\begin{zadatak}
    Sastaviti kolo neinvertujućeg pojačavača, koristiti operacioni pojačavač ADA4891, gde su $\pm V_{cc}=\pm10\,V$, a $R_1 = 100 \,k\Omega$ i $R_2 = 1 \,M\Omega$, i odrediti:
    
    \begin{enumerate}
        \item [a)] Očekivani izlaz kola ako se na ulaz dovede signal $V_{in} = 0.1\,V\,sin(20\,\pi t)$,
        \item [b)] Vrednost naponskog ofseta u kolu,
        \item [c)] Prenosnu karakteristiku pojačavača,
        \item [d)] Vrednost otpornika za kompenzaciju naponskog ofseta i
        \item [e)] Vrednost izlaznog signala za ulazni signal iz a).
    \end{enumerate}
\end{zadatak}

\pagebreak
\section{Zadatak za domaći}

\begin{zadatak}
    Za neinvertujući pojačavač realizovan preko operacionog pojačavača OP777, čija su naponi napajanja $\pm V_{cc} = \pm 10 V$, a otpornosti otpornika $R_1 = \,100\,k\Omega$ i $R_2 = \,1\,M\Omega$, odrediti:
    \begin{enumerate}
        \item [a)] Očekivani izlaz iz kola ako je ulazni signal $V_{in} = 0.1\,V \,sin(40\,\pi t)$ pod pretpostavkom idealnog pojačavača,
        \item [b)] Napisati izlaz iz kola nakon simulacije u LTSpice-u i uporediti ga sa očekivanim,
        \item [c)] Vrednost naponskog ofseta i struja polarizacije (NAPOMENA: nije moguće u isto vreme odrediti oba, prvo mora da se negira uticaj ili naponskog ofseta ili struja polarizacije kako bi se izračunala njihova zasebna vrednost),
        \item [d)] Vrednost otpornika $R_3$ za kompenzaciju struja polarizacije i vrednost otpornika $R_4$ i $R_5$ za kompenzaciju naponskog ofseta.
    \end{enumerate}
\end{zadatak}